\section{Orbit Control and Monte Carlo Simulation}
\label{sec:control}

Unstable orbits (Halo/NRHO/DPO) are only truly usable together with station
keeping. e2m2e makes control simulation the last link of the design chain:
the NominalOrbit produced during design (an equally-spaced state table +
Floquet basis + projection factors) is consumed directly by the control
algorithms, with no need to recompute the Floquet analysis.

\subsection{Three Control Laws}

\begin{itemize}
  \item \textbf{Feature-point control}: corrections are applied only at
        characteristic locations (e.g. perilune or plane crossings)---few
        maneuvers, low propellant cost, suitable for long-term keeping;
  \item \textbf{Strict target-point control}: the state is driven exactly
        back to a target point on the nominal orbit---highest tracking
        accuracy;
  \item \textbf{Relaxed target-point control}: the target point admits an
        error band, trading tracking accuracy against propellant consumption.
\end{itemize}

Angular-momentum management is also provided: joint attitude/thruster control
that folds momentum unloading into the station-keeping strategy.

\subsection{Monte Carlo Error Simulation}

Station-keeping propellant statistics must be evaluated under error models.
e2m2e supports Monte Carlo simulation with jointly injected navigation and
thrust errors: batches of keeping simulations on the same nominal orbit yield
$\Delta V$ consumption distributions and orbit-maintenance statistics, giving
mission studies confidence intervals rather than point estimates. Simulation
artifacts are automatically archived in the orbit catalog with tag
management---in the MCP scenario, ``run 100 Monte Carlo station-keeping
simulations on this NRHO and tag the results \texttt{candidate}'' is a
workflow completed by a single natural-language instruction.
